\clearpage
\section{How to Run}

From the \texttt{clustered-smart-home-alarm-system} project directory, which
contains \texttt{pom.xml}:

\begin{itemize}
  \item Build: \texttt{mvn -f pom.xml compile}
  \item Test: \texttt{mvn -f pom.xml test}
  \item Distributed demo: \texttt{mvn -f pom.xml exec:java} with main class \texttt{pcd.shas.DemoMain}
  \item Cluster node: \texttt{mvn -f pom.xml exec:java} with main class \texttt{pcd.shas.Main}
\end{itemize}

The demo entry point is implemented in
\texttt{src/main/java/pcd/shas/DemoMain.java}; the clustered-node entry point
is in \texttt{src/main/java/pcd/shas/Main.java}.

The clustered entry point \texttt{pcd.shas.Main} starts one node at a time.
The accepted roles are \texttt{control-unit}, \texttt{keypad}, and
\texttt{sensor}. A sensor node also requires \texttt{sensor-id},
\texttt{sensor-type}, and \texttt{zone}.

\textbf{Distributed demo}
\begin{lstlisting}[style=terminalstyle]
mvn -f pom.xml compile exec:java "-Dexec.mainClass=pcd.shas.DemoMain"
\end{lstlisting}

\textbf{Control-unit node}
\begin{lstlisting}[style=terminalstyle]
mvn -f pom.xml compile exec:java "-Dexec.mainClass=pcd.shas.Main" "-Dexec.args=control-unit --host 127.0.0.1 --port 2551 --seed-nodes 127.0.0.1:2551,127.0.0.1:2552,127.0.0.1:2553"
\end{lstlisting}

\textbf{Keypad node}
\begin{lstlisting}[style=terminalstyle]
mvn -f pom.xml compile exec:java "-Dexec.mainClass=pcd.shas.Main" "-Dexec.args=keypad --host 127.0.0.1 --port 2552 --seed-nodes 127.0.0.1:2551,127.0.0.1:2552,127.0.0.1:2553"
\end{lstlisting}

\textbf{Sensor node}
\begin{lstlisting}[style=terminalstyle]
mvn -f pom.xml compile exec:java "-Dexec.mainClass=pcd.shas.Main" "-Dexec.args=sensor --host 127.0.0.1 --port 2553 --sensor-id front_door --sensor-type DOOR_WINDOW --zone PERIMETER --seed-nodes 127.0.0.1:2551,127.0.0.1:2552,127.0.0.1:2553"
\end{lstlisting}

The control-unit CLI accepts \texttt{status}, \texttt{help}, and \texttt{exit}.
The keypad CLI accepts \texttt{pin PIN}, \texttt{arm full PIN},
\texttt{arm partial PIN ZONE...}, raw digit input, and \texttt{exit}. The
sensor CLI triggers one activation on ENTER and stops on \texttt{exit}.

\textbf{Example commands after startup}

\begin{itemize}
  \item Control-unit node: \texttt{status}, \texttt{help}, \texttt{exit}
  \item Keypad node: \texttt{pin 1234}, \texttt{arm full 1234},
  \texttt{arm partial 1234 PERIMETER GROUND\_FLOOR}, \texttt{1234},
  \texttt{exit}
  \item Sensor node: press ENTER to trigger one activation, then \texttt{exit}
\end{itemize}

To observe a minimal three-node execution, start one control-unit node, one
keypad node, and one sensor node with the commands above. Then submit a
correct PIN from the keypad, arm the system, wait for the exit delay to
expire, trigger the sensor from its CLI, and then submit the correct PIN
during the entry delay. The expected state evolution is
\texttt{STARTUP\_RECOVERY $\rightarrow$ DISARMED $\rightarrow$ EXIT\_DELAY
$\rightarrow$ ARMED $\rightarrow$ ENTRY\_DELAY $\rightarrow$ DISARMED}. If no
correct PIN is entered during \texttt{ENTRY\_DELAY}, the system transitions to
\texttt{ALARM}.
